% Evaluation-section text for inclusion in the team's existing LaTeX report.
% Source: our approach/evaluation.ipynb, based on main a34f744.
% This is a section fragment, not a standalone report or a new experiment.

\subsection{Evaluation protocol and metrics}
We retain the official participant-separated training and test splits, with
7,767 training windows and 3,162 test windows. Transitions account for only
352 training windows (4.53\%) and 166 test windows (5.25\%). Accordingly,
our primary metric is macro-F1 over the six transition classes (IDs 7--12),
which gives each transition equal weight. Transition macro-recall, common-class
macro-F1, overall macro-F1 and accuracy provide complementary measures.
Model development uses three-fold participant-grouped validation within the
training participants. The final test comparison describes the retained
configurations; the test participants were inspected during development and
therefore do not provide an untouched final estimate.

\subsection{Controlled comparisons}
For the original five model families, mean validation transition F1 increases
from 0.7080 with a single window to 0.8422 with three-window context, compared
with 0.7352 from class weighting alone. At fixed context, weighting improves
four of five models, with small single-window degradation for logistic
regression and a 3.20-percentage-point reduction for the three-window MLP.
Thus, context gives the clearest consistent improvement, while weighting is
model dependent. For LDA with equal priors and memory off, transition F1
increases from 0.7567 with the current window to 0.7811 with the previous
and current windows, and 0.8917 with previous/current/next windows.
LDA has no sample-weighting arm, so its missing weighting results are not
zero scores.

Random oversampling is not consistently beneficial. For example, LDA's
three-window validation transition F1 decreases to 0.8894 and 0.8881 at
targets 60 and 100, respectively. PCA also has model-specific trade-offs:
LR PCA-400 gives repeatable runtime savings within the chosen quality
tolerances, whereas LDA PCA-500 loses 1.83 percentage points of transition
F1 and 1.75 points of common F1. The retained LDA therefore uses raw features.

\subsection{Comparison of final configurations}
Table~\ref{tab:final-evaluation} compares all six retained pipelines. The first
five models use mild class weighting; LDA uses equal priors. MLP and LDA
include posture-memory postprocessing. No final configuration uses oversampling.
Different window lengths and class treatments mean this comparison does not
isolate the effect of classifier family.

\begin{table*}[t]
\centering
\small
\caption{Recorded scores for the retained configurations. CV is the mean
participant-grouped validation transition F1; the other columns use the official
test windows. Only LR uses PCA-400.}
\label{tab:final-evaluation}
\begin{tabular}{lrrrrrr}
\hline
Model & Windows & CV & Trans. F1 & Common F1 & Macro F1 & Accuracy \\
\hline
Logistic regression & 7 & 0.9048 & 0.9373 & 0.9615 & 0.9494 & 0.9592 \\
Linear SVM          & 7 & 0.8834 & 0.9322 & 0.9784 & 0.9553 & 0.9741 \\
Random forest       & 9 & 0.9201 & 0.9296 & 0.9344 & 0.9320 & 0.9345 \\
XGBoost             & 7 & 0.9033 & 0.9285 & 0.9696 & 0.9491 & 0.9681 \\
MLP + memory        & 11 & 0.9302 & 0.9716 & 0.9766 & 0.9741 & 0.9756 \\
LDA + memory        & 7 & 0.9699 & 0.9779 & 0.9787 & 0.9783 & 0.9782 \\
\hline
\end{tabular}
\end{table*}

Seven-window shrinkage LDA with posture memory has the highest validation
transition F1 among the retained configurations (0.9699). Its test transition
F1 is 0.9779, compared with 0.9716 for the eleven-window MLP, a difference
of 0.63 percentage points using the rounded scores. LDA's transition
macro-precision and macro-recall are 0.9829 and 0.9736. These descriptive
differences do not establish statistical significance or general superiority.

\subsection{Error analysis and limitations}
Recomputing LDA's metrics from all 3,162 saved test predictions reproduces
the recorded confusion matrix and aggregate scores. LDA correctly labels
161 of 166 true transition windows, versus MLP's recorded 158. LDA's
\texttt{STAND\_TO\_SIT} has the lowest transition F1 (0.9565), while
\texttt{LIE\_TO\_STAND} has the lowest transition recall (25/27, or 92.59\%).
The five true-transition errors comprise three predictions as common activities
and two predictions as another transition. Four belong to participant 24.
One additional common window is falsely predicted as a transition:
\texttt{LAYING} to \texttt{STAND\_TO\_LIE}.
Of LDA's 69 total errors, 48 are sitting/standing swaps and 13 are
walking-downstairs windows predicted as another walking activity.

Full MLP per-window predictions are not archived in the current source
notebooks, so its full confusion matrix, per-class precision and error
destinations cannot be reconstructed from correct/support counts alone.
Nor do F1-only weighting tables establish matched recall or common-class
trade-offs. Future-window context describes offline inference, and the prepared
CSVs lack timestamps and within-participant recording IDs. Window-level results
therefore do not establish online or event-level performance. Recording-aware
evaluation and preserving complete predictions would address these limitations;
the observed test errors should not be used to retune the selected pipelines.
